\documentclass[aspectratio=169,10pt,english]{beamer}

\input{../../../_preambulo_beamer.tex}
\input{../../../_comandos_beamer.tex}

\selectlanguage{english}

\title{MA1001B - Statistical Modeling for Decision Making}
\subtitle{Lesson 23: Sampling Distribution of the Mean}
\author{}
\institute{Tecnol\'ogico de Monterrey}
\date{}

\begin{document}

\begin{frame}[plain]
  \vspace{1.2cm}
  {\Large\bfseries MA1001B - Statistical Modeling for Decision Making\par}
  \vspace{0.7cm}
  {\large Lesson 23: Sampling Distribution of the Mean\par}
  \vspace{0.7cm}
  {\color{TecRojo}\rule{\textwidth}{1.2pt}}
  \vspace{0.8cm}

  MA1001B Faculty\\[0.3cm]
  Term\\[0.3cm]
  Tecnol\'ogico de Monterrey
\end{frame}

\begin{frame}{The Sample-Mean Question}
  \begin{alertblock}{Guiding question}
    If cycle times are $N(80,12)$, how much does $\bar x$ move when
    $n=9$ versus $n=36$?
  \end{alertblock}

  \vspace{0.6em}
  By the end of this lesson, you can:
  \begin{itemize}
    \item write $E[\bar x]=\mu$ and $\mathrm{SE}(\bar x)=\sigma/\sqrt{n}$;
    \item compare sampling distributions for two sample sizes;
    \item check the SE formula with a seeded simulation;
    \item explain why clustered copies inflate the observed SE.
  \end{itemize}

  \vfill
  \scriptsize All cycle times used today are synthetic. No real company data are used.
\end{frame}

\begin{frame}{Standard Error Shrinks with $n$}
  \small
  Population $X\sim N(80,12^2)$. For an iid sample,
  \[
    E[\bar x]=80,\qquad
    \mathrm{SE}(\bar x)=\frac{12}{\sqrt{n}}.
  \]
  Then
  \[
    \mathrm{SE}(9)=\frac{12}{3}=4,\qquad
    \mathrm{SE}(36)=\frac{12}{6}=2.
  \]
  Quadrupling $n$ halves the SE.

  \vspace{0.3em}
  \begin{exampleblock}{Synthetic cycle-time population}
    $\sigma=12$ describes one observation. $\sigma/\sqrt{n}$ describes $\bar x$.
  \end{exampleblock}
\end{frame}

\begin{frame}[fragile]{Step 1: Standard Error Formula}
  \begin{columns}[T]
    \begin{column}{0.42\textwidth}
      \small
      \begin{lstlisting}
se_9 = 12 / sqrt(9)
se_36 = 12 / sqrt(36)
      \end{lstlisting}

      \begin{block}{Verified output}
        SE$(n=9)=4.000000$\\
        SE$(n=36)=2.000000$\\
        Ratio $=2.000000$
      \end{block}
    \end{column}
    \begin{column}{0.58\textwidth}
      \centering
      \includegraphics[width=\textwidth]{../figures/sampling_mean_01_standard_error.png}
      \vspace{0.2em}
      \scriptsize Population versus sampling densities from Step 1
    \end{column}
  \end{columns}
\end{frame}

\begin{frame}{Simulation Recovers the SE Formula}
  \small
  Draw 5{,}000 iid samples with seed 42.
  \[
    n=9:\ \text{mean of means}=80.036617,\ \widehat{\mathrm{SE}}=3.997160.
  \]
  \[
    n=36:\ \text{mean of means}=79.994341,\ \widehat{\mathrm{SE}}=2.028859.
  \]
  Both empirical SEs sit next to the formulas 4 and 2.
\end{frame}

\begin{frame}[fragile]{Step 2: $n=9$ versus $n=36$}
  \begin{columns}[T]
    \begin{column}{0.40\textwidth}
      \small
      \begin{lstlisting}
draws = rng.normal(
  80, 12, size=(5000, n)
)
xbars = draws.mean(axis=1)
      \end{lstlisting}

      \begin{block}{Verified output}
        $n=9$ SE $=3.997160$\\
        $n=36$ SE $=2.028859$
      \end{block}
    \end{column}
    \begin{column}{0.60\textwidth}
      \centering
      \includegraphics[width=\textwidth]{../figures/sampling_mean_02_n9_vs_n36.png}
      \vspace{0.2em}
      \scriptsize Simulated sampling distributions from Step 2
    \end{column}
  \end{columns}
\end{frame}

\begin{frame}{Step 3: SE Assumes iid Draws}
  \begin{columns}[T]
    \begin{column}{0.42\textwidth}
      \small
      Each ``sample'' of 36 is 6 cluster values copied 6 times.

      \vspace{0.4em}
      Formula SE for $n=36$ is $2$.\\
      Clustered empirical SE $=4.992583$.\\
      Formula SE for 6 clusters is $4.898979$.

      \vspace{0.4em}
      \textbf{Limit:} SE assumes iid draws from the defined population.
      Lesson 24 invokes the CLT for a skewed clock.
    \end{column}
    \begin{column}{0.58\textwidth}
      \centering
      \includegraphics[width=\textwidth]{../figures/sampling_mean_03_iid_limit.png}
      \vspace{0.2em}
      \scriptsize iid versus clustered means from Step 3
    \end{column}
  \end{columns}
\end{frame}

\begin{frame}{Concept Check}
  \begin{enumerate}
    \item Compute $\mathrm{SE}(\bar x)$ for $n=9$ when $\sigma=12$.
    \item Why does quadrupling $n$ from 9 to 36 halve the SE?
    \item The $n=36$ empirical SE is $2.028859$. Does that replace $2$?
    \item Why is a sample of 36 clustered copies closer to $n=6$?
  \end{enumerate}

  \vspace{0.6em}
  \begin{block}{Connection to Lesson 24}
    The session can continue directly into the central limit theorem for a
    right-skewed population.
  \end{block}

  \vfill
  \scriptsize Source: OpenStax, \textit{Introductory Business Statistics 2e},
  Chapter 7, Section 7.1. CC BY-NC-SA.
\end{frame}

\appendix



\end{document}
